\section{Modelo de Implementación}

Este apartado documenta el entorno tecnológico sobre el que se construyó el
sistema, la forma en que sus componentes se despliegan y la organización del
código fuente. La descripción es deliberadamente concreta: frente al capítulo de
diseño, que establece la arquitectura conceptual, aquí se registran las
tecnologías efectivamente empleadas y su disposición física, de modo que la
implementación pueda reproducirse sobre otro entorno.

\subsection{Entorno tecnológico y dependencias}

La selección de tecnologías responde a las restricciones identificadas en el
diseño: la necesidad de ejecutar inferencia profunda sobre hardware con recursos
limitados, de desacoplar el trabajo pesado de la atención de solicitudes y de
mantener una única fuente de verdad para los datos. La Tabla
\ref{tab:entorno} resume el entorno tecnológico.

\begin{longtable}{
        P{0.18}
        P{0.35}
        P{0.39}
    }
    \caption{Entorno tecnológico del sistema}
    \label{tab:entorno}\\

    \toprule
    \rowcolor{gray!20}
    \textbf{Componente} &
    \textbf{Tecnología y versión} &
    \textbf{Rol en el sistema} \\
    \midrule
    \endfirsthead

    \multicolumn{3}{c}{\textit{Continuación de la Tabla
    \ref{tab:entorno}}} \\
    \toprule
    \rowcolor{gray!20}
    \textbf{Componente} &
    \textbf{Tecnología y versión} &
    \textbf{Rol en el sistema} \\
    \midrule
    \endhead

    Lenguaje y entorno & Python 3.12; gestor de dependencias \textit{uv} con
    archivo de bloqueo & Servicios, trabajador y herramientas de datos \\
    \midrule

    Servicio web & FastAPI $\geq 0{.}115$; Uvicorn $\geq 0{.}32$ & API REST y
    documentación automática de la interfaz \\
    \midrule

    Acceso a datos & SQLAlchemy $\geq 2{.}0$ en modo asíncrono; asyncpg
    $\geq 0{.}30$ & Persistencia no bloqueante desde el servicio \\
    \midrule

    Migraciones & Alembic $\geq 1{.}14$ & Versionado del esquema de datos \\
    \midrule

    Base de datos & PostgreSQL 16 & Almacenamiento relacional de vehículos,
    sesiones y detecciones \\
    \midrule

    Cola de tareas & Celery $\geq 5{.}4$; Redis $\geq 5{.}1$ & Ejecución
    asíncrona de la inferencia \\
    \midrule

    Detección de objetos & Ultralytics (YOLO26) $\geq 8{.}4$ & Localización de
    vehículos y placas de matrícula \\
    \midrule

    Aprendizaje profundo & PyTorch $\geq 2{.}4$; torchvision $\geq 0{.}19$ &
    Entrenamiento e inferencia de los modelos \\
    \midrule

    Visión clásica & OpenCV (sin interfaz gráfica) $\geq 4{.}10$ & Captura,
    preprocesamiento y calibración de zonas \\
    \midrule

    Seguimiento & ByteTrack mediante \textit{lapx} $\geq 0{.}5{.}2$ & Asociación
    temporal de detecciones entre fotogramas \\
    \midrule

    Reconocimiento de texto & PaddleOCR $\geq 2{.}8$; PaddlePaddle $\geq 3{.}0$ &
    Pseudo-etiquetado de placas y comparación de referencia \\
    \midrule

    Optimización de modelos & ONNX, onnxslim, ONNX Runtime GPU y TensorRT &
    Exportación y aceleración de la inferencia \\
    \midrule

    Notificación & \textit{websockets} $\geq 13$ & Canal persistente con el panel
    de monitoreo \\
    \midrule

    Interfaz web & React 19; Vite 6; Tailwind CSS 3.4 & Aplicación del personal
    del parqueadero \\
    \midrule

    Estado del cliente & TanStack Query 5; TypeScript 5.7 & Sincronización con la
    API y tipado del cliente \\
    \midrule

    Contenerización & Docker; Docker Compose & Despliegue reproducible de los
    seis contenedores \\*
    \multicolumn{3}{r}{\footnotesize\textit{Fuente: Elaboración propia}} \\*
    \bottomrule
\end{longtable}

La gestión de dependencias se realiza con un archivo de bloqueo que fija las
versiones exactas de cada paquete, de modo que la instalación resulte
reproducible. La resolución se restringe explícitamente a Python 3.12, dado que
las versiones posteriores del intérprete imponen una versión de la biblioteca
numérica incompatible con la requerida por el motor de reconocimiento de texto.

\subsection{Arquitectura de despliegue}

El sistema se despliega como un conjunto de seis contenedores que
comparten una red interna y se comunican por nombre de servicio: cuatro servicios
de aplicación ---API, trabajador, planificador y cliente web---, una base de datos
y un intermediario de cola. La Figura \ref{fig:despliegue} presenta su
disposición; la Tabla \ref{tab:servicios} detalla la función de cada uno y sus
dependencias de arranque. El planificador se omite de la figura por claridad: su
única función es disparar la rutina diaria de retención.

\begin{figure}[htbp]
    \centering
    \begin{tikzpicture}
        \spBlock{frontend}{-5.0,0}{3.6cm}{cielo}
            {Cliente web}{Nginx}
        \spBlock{api}{0,0}{3.6cm}{cielo}
            {API REST}{FastAPI + Uvicorn}
        \spBlock{bd}{5.0,0}{3.6cm}{purpura}
            {PostgreSQL}{Datos del sistema}

        \spBlock{redis}{0,-2.8}{3.6cm}{purpura}
            {Redis}{Cola de mensajes}
        \spBlock{worker}{5.0,-2.8}{3.6cm}{cielo}
            {Worker}{Celery + inferencia}

        \draw[spflow] (frontend.east) --
            node[above, font=\tiny\itshape, text=black!70] {HTTP}
            (api.west);
        \draw[spflow] (api.east) --
            node[above, font=\tiny\itshape, text=black!70] {SQL}
            (bd.west);
        \draw[spflow] (api.south) -- (redis.north);
        \draw[spflow] (worker.west) -- (redis.east);
        \draw[spflow] (worker.north) -- (bd.south);

    \end{tikzpicture}
    \caption{Arquitectura de despliegue en contenedores.}
    \label{fig:despliegue}
\end{figure}

\begin{longtable}{
        P{0.10}
        P{0.28}
        P{0.54}
    }
    \caption{Servicios de despliegue y sus dependencias}
    \label{tab:servicios}\\

    \toprule
    \rowcolor{gray!20}
    \textbf{Servicio} &
    \textbf{Tecnología} &
    \textbf{Función y dependencias} \\
    \midrule
    \endfirsthead

    \multicolumn{3}{c}{\textit{Continuación de la Tabla
    \ref{tab:servicios}}} \\
    \toprule
    \rowcolor{gray!20}
    \textbf{Servicio} &
    \textbf{Tecnología} &
    \textbf{Función y dependencias} \\
    \midrule
    \endhead

    \texttt{api} & FastAPI sobre Uvicorn & Expone los recursos REST y el canal de
    notificación. Depende de que la base de datos esté disponible y de que la cola
    se encuentre en ejecución. \\
    \midrule

    \texttt{worker} & Celery & Consume las tareas de inferencia encoladas y
    actualiza los resultados en la base de datos. Comparte las mismas
    dependencias de arranque que el servicio. \\
    \midrule

    \texttt{beat} & Celery Beat & Ejecuta la rutina diaria de retención de
    telemetría y recortes. Depende de la base de datos y de la cola; su caída no
    detiene la inferencia. \\
    \midrule

    \texttt{postgres} & PostgreSQL 16 & Almacena el estado del sistema sobre un
    volumen persistente. Declara una comprobación de salud que condiciona el
    arranque de los servicios que la consumen. \\
    \midrule

    \texttt{redis} & Redis 7 & Actúa como intermediario de mensajes y como
    almacén de resultados de las tareas. \\
    \midrule

    \texttt{frontend} & React servido por Nginx & Entrega la aplicación web
    compilada y depende del servicio de la API. \\*
    \multicolumn{3}{r}{\footnotesize\textit{Fuente: Elaboración propia}} \\*
    \bottomrule
\end{longtable}

Los servicios de aplicación se construyen a partir de imágenes base de Python
3.12 en dos etapas: una de compilación, en la que se resuelven las dependencias
mediante el gestor de paquetes, y otra de ejecución, que recibe únicamente lo
estrictamente necesario. El cliente web se compila en una imagen de Node y se
sirve como contenido estático sobre Nginx. Los datos se conservan en volúmenes
declarados explícitamente, de modo que la reconstrucción de los contenedores no
suponga pérdida de información.

El trabajador de inferencia se ejecuta con una única unidad de concurrencia, ya
que cada tarea satura la unidad de procesamiento gráfico disponible. La
reservación del dispositivo gráfico para este contenedor se declara en el archivo
de composición y se habilita según el entorno de despliegue: en una estación de
trabajo con acelerador dedicado se activa, mientras que en un entorno sin
aceleración el trabajador opera sobre el procesador.

\subsection{Estructura del repositorio y organización del código}

El código se organiza en cuatro áreas: el servicio de gestión, el trabajador de
inferencia, el cliente web y las herramientas de datos y de despliegue de
modelos. La Tabla \ref{tab:repositorio} describe la responsabilidad de cada
directorio.

\begin{longtable}{
        P{0.33}
        P{0.62}
    }
    \caption{Organización del código fuente}
    \label{tab:repositorio}\\

    \toprule
    \rowcolor{gray!20}
    \textbf{Ruta} &
    \textbf{Responsabilidad} \\
    \midrule
    \endfirsthead

    \multicolumn{2}{c}{\textit{Continuación de la Tabla
    \ref{tab:repositorio}}} \\
    \toprule
    \rowcolor{gray!20}
    \textbf{Ruta} &
    \textbf{Responsabilidad} \\
    \midrule
    \endhead

    \texttt{backend/app/api/routes/} & Definición de los recursos HTTP y del canal
    de notificación. \\
    \midrule

    \texttt{backend/app/models/} & Modelos de datos: vehículo, sesión de
    estacionamiento y detección. \\
    \midrule

    \texttt{backend/app/schemas/} & Contratos de entrada y salida de la interfaz
    de programación. \\
    \midrule

    \texttt{backend/app/core/} & Configuración, sesión de base de datos y cliente
    de la cola. \\
    \midrule

    \texttt{backend/app/main.py} & Composición de la aplicación y gestión del
    ciclo de vida. \\
    \midrule

    \texttt{backend/worker/} & Aplicación Celery y tareas de inferencia. \\
    \midrule

    \texttt{backend/alembic/} & Migraciones versionadas del esquema. \\
    \midrule

    \texttt{backend/tests/} & Pruebas automatizadas del servicio. \\
    \midrule

    \texttt{frontend/src/pages/} & Vistas del panel de monitoreo y del historial. \\
    \midrule

    \texttt{frontend/src/components/} & Componentes reutilizables de la interfaz. \\
    \midrule

    \texttt{frontend/src/hooks/} & Enlaces del cliente, incluido el canal
    persistente con el servicio. \\
    \midrule

    \texttt{frontend/src/services/} & Acceso a la API y gestión del canal de
    notificación. \\
    \midrule

    \texttt{rtsp-capture/} & Captura de flujos RTSP y despacho de tareas de
    inferencia. \\
    \midrule

    \texttt{tools/} & Calibración de zonas, extracción del conjunto de datos y
    exportación de modelos. \\*
    \multicolumn{2}{r}{\footnotesize\textit{Fuente: Elaboración propia}} \\*
    \bottomrule
\end{longtable}

La separación entre el servicio de gestión y el trabajador de inferencia se
refleja en el propio árbol de directorios: el primero concentra las reglas de
negocio y el acceso a los datos, mientras que el segundo aloja exclusivamente la
lógica de procesamiento. Las herramientas del directorio \texttt{tools} se
ejecutan fuera del ciclo de operación y no forman parte de los servicios
desplegados.
